\section{Einleitung}
\epigraph{"`Des Menschen größtes Verdienst bleibt wohl, wenn er die Umstände soviel als möglich bestimmt und sich so wenig als möglich von ihnen bestimmen läßt."'}{Johann Wolfgang von Goethe\footcite[S. 10]{Freund2014}}
Ein zum Thema passendes Zitat fast immer eine gute Einleitung für die Arbeit. 
\ac{BPMN} ist eine Modellierungssprache.\footcite[Vgl.][S. 1]{Freund2014} Bei der ersten Verwendung von Abkürzungen werden diese in Klammern automatisch ausgeschrieben. 
Bei der zweiten Verwendung ist das nicht so, wie man anhand von \ac{BPMN} sehen kann. "`Das ist ein direktes Zitat aus dem Internet"'.\footcite[S. 3]{OMG2018}
Bei einseitigen Quellen kann man die Seitenzahl weglassen.\footcite[Vgl.][]{schlechteQuelle}
Es gibt viele schlechte Quellen.\footcite[Vgl.][S. 1-3]{schlechteQuelle2}

\subsection{Motivation}
Abbildungen und Tabellen sind natürlich auch möglich.

\begin{figure}[H]
  \centering
	\includegraphics[width=0.4\textwidth]{DHBW_logo.pdf}
   \caption[Das Logo der DHBW]{Das Logo der DHBW\footnotemark}
\end{figure}
\footnotetext{\cite[][S.1]{DHBW}}

Als Grafikformate werden u.a. PDF, PNG und JPEG akzeptiert. Die Bilddatei muss im Order "`images"' liegen.
Mit einem Label in einer Abbildung oder Tabelle kann man darauf referenzieren, wie man an der Abbildung \ref{AbbildungLogoMusterfirma} sehen kann.
\begin{figure}[H]
  \centering
	\includegraphics[width=0.3\textwidth]{company_logo.pdf}
   \caption[Das Logo der Musterfirma]{Das Logo der Musterfirma\protect\footnotemark}
   \label{AbbildungLogoMusterfirma}
\end{figure}
\footnotetext{Eigene Darstellung in Anlehnung an \cite[S.4]{Freund2014}}

Die Breite einer Grafik oder einer Tabelle lässt sich einfach als Faktor festlegen. 1 entspricht dabei der Textbreite und 0.5 die Hälfte der Textbreite.
Bei Tabellen wird die angegebene Breite nur bei Bedarf ausgenutzt.

\begin{table}[H]
  \centering
  \begin{tabulary}{0.7\textwidth}{|L|L|}
  \hline 
  \rowcolor{tableHeading}Eigenschaft & Wert \\ 
  \hline 
  Größe & 20 cm \\ 
  \hline 
  Gewicht & 1 kg \\
  \hline
  Haarfarbe & braun \\  
  \hline 
  \end{tabulary} 
  \caption[Eine Tabelle ohne Quellenangabe]{Eine Tabelle ohne Quellenangabe}
\end{table}


Experteninterviews.\footcite[Vgl.][S. 7]{Meuser2009} Ein Zitat aus der Wirtschaftswoche.\footcite[Vgl.][S. 32]{WIWO2018}
Firmeninternes Material kann auch zitiert werden.\footcite[Vgl.][]{Firma2018} "`Das Zitat stammt aus einem Interviewprotokoll"'.\footcite{Experte2018}

Mit zwei Backslash \\ erzwingt man einen Zeilenumbruch. Bei langen Wörtern funktioniert die Worttrennung oftmals nicht mehr.
Dann muss man selbst die Silbentrennung vornehmen: Donau\-dampf\-schiff\-fahrts\-gesell\-schafts\-kapitän. \\
Aufzählungen:
\begin{itemize}
  \item Punkt 1
  \item Punkt 2
\end{itemize}
Nummerierte Aufzählung:
\begin{enumerate}
  \item Punkt 1
  \item Punkt 2
\end{enumerate}

Fußnoten sind besonders praktisch für Verweise auf andere Abschnitte der Arbeit.\footnote{Siehe Abschnitt \ref{Abschnitt:Arbeitsumfeld}} 
Mit dem ref-Befehl lassen sich Labels referenzieren. Das funktioniert bei Abbildungen, Tabellen, Kapiteln und Abschnitten.